\section{Bevisindex}

\subsection{Inledande matematisk analys}

\subsubsection{Faktorsatsen för polynom}
\begin{theorem}[Faktorsatsen för polynom]
	$(x-a)$ är en faktor till polynomet $p\iff p(a)=0$.

	\begin{proof}
		Med polynomdivision av $p$ genom $(x-a)$ kan vi skriva $p(x)$ som:
		$$
			p(x)=(x-a)q(x)+r(x)
		$$
		där det gäller att $\pgrad r<\pgrad (x-a)=1\implies r$ är en konstant.
		Eftersom $a$ är ett nollställe till $p$, så vet vi att:
		\begin{align*}
			0&=p(a)=(a-a)q(a)+r\\
			&=r
		\end{align*}
		Därmed kan $p$ skrivas enligt:
		$$
			p(x)=(x-a)q(x)
		$$
		omm $a$ är ett nollställe till $p$.
	\end{proof}
\end{theorem}

\subsubsection{Eventuella rationella nollställen till polynom}
\begin{theorem}[Eventuella rationella nollställen till polynom]
	Satsen ger en begräning till lösningarna av polynom ekvationen $\sum_{k=0}^{n}c_kx^k=0$,
	där $c_k\in\mathbb{Z}$ och $c_0,c_n\neq0$. Alla lösningar $r=\frac{a}{b}\in\mathbb{Q}$
	där $\gcd(a,b)=1$ uppfyller:
	\begin{itemize}
		\item $\gcd(a,a_0)\neq1$
		\item $\gcd(b,a_n)\neq1$
	\end{itemize}

	\begin{proof}
		Låt $p:x\mapsto\sum_{k=0}^{n}c_kx^k=0$, där $c_k\in\mathbb{Z}$ och $c_0,c_n\neq0$.
		Antag $p(\frac{a}{b})=0$ för något par $a,b\in\mathbb{Z}$, där $\gcd(a,b)=1$.
		Detta ger:
		$$
			p\qty(\frac{a}{b})=\sum_{k=0}^{n}a_k\qty(\frac{a}{b})^k=0
		$$
		Genom multiplikation med $q^n$, ges:
		\begin{align*}
			b^np\qty(\frac{a}{b})=c_na^n+c_{n-1}a^{n-1}b+\dots+c_0b^n&=0\\
			\implies a(c_na^{n-1}+c_{n-1}a^{n-2}b+\dots+c_1b^{n-1})&=-c_0b^n
		\end{align*}
		Eftersom $c_0,b\in\mathbb{Z}$, så gäller det att $c_0b^k\equiv 0\mod a$,
		men då $a$ och $b$ är koprima tal (vilket innebär $\gcd(a,b^n)=1$), så måste
		det vara att $c_0\equiv 0\mod a$. Liknande, fast för $b$, gäller:
		\begin{align*}
			b^np\qty(\frac{a}{b})=c_na^n+c_{n-1}a^{n-1}b+\dots+c_0b^n&=0\\
			\implies b(c_{n-1}a^{n-1}+c_{n-2}a^{n-2}b+\dots+c_0b^{n-1})&=-c_na^n\\
			\therefore c_n&\equiv 0\mod b
		\end{align*}
	\end{proof}
\end{theorem}

\subsubsection{Binomialsatsen}
\begin{theorem}[Binomialsatsen]
	Låt $x,y\in\mathbb{R}$ vara sådan att $x+y\neq0$. $\forall n\in\mathbb{N}$
	gäller det att:
	$$
		(x+y)^n=\sum_{k=0}^{n}{n\choose k}x^{n-k}y^k 
	$$

	\begin{proof}
		Låt $x$ och $y$ vara godtyckliga tal enligt ovan. Eftersom
		$$
			(x+y)^n=\qty(y\qty(\frac{x}{y}+1))^n=y^n\qty(1+\frac{x}{y})^n
		$$
		så räcker det att bevisa satsen med $y=1$ u.f.a. (utan att förlora i
		allmängiltighet). Satsen bevisas genom induktion. Låt $p$ vara utsagan
		$$
			p_n:(1+x)^n=\sum_{k=0}^{n}{n\choose k}1^{n-k}x^k=\sum_{k=0}^n{n\choose k}x^k
		$$
		För $n=1$, gäller det:
		$$
			p_1:(1+x)^1=1+x={1\choose 0}1^{1}x^0+{1\choose 1}1^0x^1\quad\checkmark
		$$
		Antag sedan för något fixed $m\in\mathbb{N}\cap[2,\infty)$ att $p_m$ är
		sant. Då, för $p_{m+1}$, gäller\footnotemark:
		\begin{align*}
			(1+x)^{m+1}&=(1+x)(1+x)^m\\
			&\stackrel{p_m}{=}(1+x)\sum_{k=0}^{m}{m\choose k}x^k\\
			&=\sum_{k=0}^{m}{m\choose k}x^k+\sum_{k=0}^{m}{m\choose k}x^{k+1}\\
			&=1+\sum_{k=1}^{m}{m\choose k}x^k+\sum_{k=1}^m{m\choose k-1}x^k+x^{m+1}\\
			&=1+\qty(\sum_{k=1}^{m}\qty[{m\choose k}+{m\choose k-1}]x^k)+x^{m+1}\\
			&={m+1\choose 0}x^0+\qty(\sum_{k=1}^{m}{m+1\choose k}x^k)+{m+1\choose m+1}x^{m+1}\\
			&=\sum_{k=0}^{m+1}{m+1\choose k}x^k
		\end{align*}
		Därmed, eftersom $p_0=\top$, och $p_m\implies p_{m+1}$, för något fixed
		$m\in\mathbb{N}$, så följer det av induktionsaxiomet att $p_n$ är sant
		$\forall n\in\mathbb{N}$.
	\end{proof}

	\footnotetext{Notera $
		\sum_{k=0}^{m}{m\choose k}x^{k+1}=\sum_{k=1}^{m+1}{m\choose k-1}x^{k}=\qty[\sum_{k=1}^{m}{m\choose k-1}x^{k}]+x^{m+1}
	$}
\end{theorem}

\subsubsection{Standardgränsvärden}

\begin{theorem}[Ett standardgränsvärde]
	
\end{theorem}

\subsubsection{Gränsvärdesräkneregler}

\begin{theorem}[Sammansättning av gränsvärden]\label{sats:sam-g-v.}
	Om $\lim_{x\to B}f(x)=A$ och $\lim_{x\to a}g(x)=B$, så gäller det att:
	$$
		\lim_{x\to a}(f\circ g)(x)=A
	$$
	
	\begin{proof}
		Vi vill visa att $\forall\eps>0\exists\delta>0$ s.a.
		$$
			0<\qty|x-a|<\delta\implies\qty|(f\circ g)(x)-A|<\eps.
		$$
		Eftersom både $f$ och $g$:s gränsvärden existerar vid $a$, så kan vi välja
		$\delta_1$ och $\delta_2$ för respektive funktioner, s.a.:
		\begin{align*}
			0<\qty|y-B|<\delta_1\implies&\qty|f(y)-A|<\eps\\
			0<\qty|x-a|<\delta_2\implies&\qty|g(x)-B|<\delta_1.
		\end{align*}
    	Välj sedan $\delta=\delta_2$. Då gäller att om $0<|x-a|<\delta$, så är
		$0<|g(x)-B|<\delta_1$ vilket innebär $|f(g(x)) - A|<\eps$. Altså,
		$$
			0<|x-a|<\delta\implies|f(g(x)) - A|<\eps.
		$$
    	Därmed följer det ur definitionen av gränsvärden att
		$$
			\lim_{x \to a}(f \circ g)(x)=A.
		$$
	\end{proof}
\end{theorem}

\begin{theorem}[Addition av gränsvärden]\label{sats:add-g-v.}
	Om $\lim_{x\to a}f(x)=A$ och $\lim_{x\to a}g(x)=B$, så gäller det att:
	$$
		\lim_{x\to a}f(x)+g(x)=A+B
	$$
	
	\begin{proof}
		Vi vill visa att $\forall\eps>0\exists\delta>0$ s.a.
		$$
			0<\qty|x-a|<\delta\implies\qty|f(x)+g(x)-(A+B)|<\eps.
		$$
		Eftersom både $f$ och $g$:s gränsvärden existerar vid $a$, så kan vi välja
		$\delta_1$ och $\delta_2$ för respektive funktioner, s.a.:
		\begin{align*}
			0<\qty|x-a|<\delta_1\implies&\qty|f(x)-A|<\eps_1\leteq\frac{\eps}{2}\\
			0<\qty|x-a|<\delta_2\implies&\qty|g(x)-B|<\eps_2\leteq\frac{\eps}{2}
		\end{align*}
		Låt sedan $\delta=\min\qty{\delta_1,\delta_2}$\footnotemark. Då följer det
		att:
		\begin{align*}
			0<\qty|x-a|<\delta\implies\qty|f(x)-A+g(x)-B|
			&\leq\qty|f(x)-A|+\qty|g(x)-B|\\
			&<\eps_1+\eps_2\\
			&<\frac{\eps}{2}+\frac{\eps}{2}=\eps
		\end{align*}
		Altså följer det av definitionen av gränsvärden att:
		$$
			\lim_{x\to a}f(x)+g(x)=A+B
		$$
	\end{proof}

	\footnotetext{För det gränsvärdet
		med det kanoniskt större $\delta$-värde, gäller fortfarande följden då
		$(a-\delta_m,a+\delta_m) \subset(a-\delta_M,a+\delta_M)$, där
		$\delta_m\leteq\min\qty{\delta_1,\delta_2}$ och
		$\delta_M\leteq\max\qty{\delta_1,\delta_2}$.}
\end{theorem}

\begin{theorem}[Multiplikation av gränsvärden]\label{sats:mult-g-v.}
	Om $\lim_{x\to a}f(x)=A$ och $\lim_{x\to a}g(x)=B$, så gäller det att:
	$$
		\lim_{x\to a}f(x)g(x)=AB
	$$
	
	\begin{proof}
		Låt $\eps>0$. Det existerar då $\delta_1,\delta_2>0$ s.a.:
		\begin{align*}
			0<\qty|x-a|<\delta_1\implies&\qty|f(x)-A|<\sqrt{\eps}\\
			0<\qty|x-a|<\delta_2\implies&\qty|g(x)-B|<\sqrt{\eps}
		\end{align*}
		Definera sedan $\delta=\min\qty{\delta_1,\delta_2}$. $\forall0<\qty|x-a|<\delta$
		får vi då:
		\begin{align*}
			\qty|(f(x)-A)(g(x)-B)-0|
			&=\qty|f(x)-A|\qty|g(x)-B|\\
			&<\sqrt{\eps}\sqrt{\eps}\\
			&=\eps.
		\end{align*}
		Altså är
		$$
			\lim_{x\to a}(f(x)-A)(g(x)-B)=0.
		$$
		Vi distruberar multiplikationen:
		\begin{align*}
			(f(x)-A)(g(x)-B)&=f(x)g(x)-Ag(x)-Bf(x)+AB\\
			\implies f(x)g(x)&=(f(x)-A)(g(x)-B)+Ag(x)+Bf(x)-AB.
		\end{align*}
		Genom att applicera Sats~\ref{sats:add-g-v.} kan vi bestäma:
		\begin{align*}
			\lim_{x\to a}f(x)g(x)&=\lim_{x\to a}\qty[(f(x)-A)(g(x)-B)+Ag(x)+Bf(x)-AB]\\
			&=\lim_{x\to a}(f(x)-A)(g(x)-B)+\lim_{x\to a}Ag(x)+\lim_{x\to a}Bf(x)-\lim_{x\to a}AB\\
			&=0+AB+BA-AB\\
			&=AB
		\end{align*}
	\end{proof}
\end{theorem}

\begin{theorem}[Division av gränsvärden]
	Om $\lim_{x\to a}f(x)=A$ och $\lim_{x\to a}g(x)=B$, så gäller det att:
	$$
		\lim_{x\to a}\frac{f(x)}{g(x)}=\frac{A}{B}
	$$
	
	\begin{proof}
		Låt $\eps>0$. Eftersom $\lim_{x\to a}g(x)=B$ existerar det något $\delta_1>0$
		s.a.:
		$$
			0<\qty|x-a|<\delta_1\implies\qty|g(x)-B|<\frac{\qty|B|}{2}
		$$
		Antagande $0<\qty|x-a|<\delta_1$ gäller det att:
		\begin{align}
			\notag\qty|B|
			\notag&=\qty|B-g(x)+g(x)|\\
			\notag&\leq\qty|B-g(x)|+\qty|g(x)|\\
			\notag&=\qty|B-g(x)|+\qty|g(x)|\\
			\notag&<\frac{\qty|B|}{2}+\qty|g(x)|\\
			\notag\implies\frac{\qty|B|}{2}&<\qty|g(x)|\\
			\label{eq:sats:div-g-v.:1}\implies\frac{1}{\qty|g(x)|}&<\frac{2}{\qty|B|}
		\end{align}
		Nu, låt det även finnas ett $\delta_2>0$ s.a.:
		\begin{equation}
			\label{eq:sats:div-g-v.:2}0<\qty|x-a|<\delta_2\implies\qty|g(x)-B|<\frac{\qty|B|^2}{2}\eps
		\end{equation}
		Vi väljer sedan $\delta=\min\qty{\delta_1,\delta_2}$. Om $0<\qty|x-a|<\delta$
		har vi:
		\begin{align*}
			\qty|\frac{1}{g(x)}-\frac{1}{B}|
			&=\qty|\frac{B-g(x)}{Bg(x)}|\\
			&=\frac{1}{\qty|Bg(x)|}\qty|B-g(x)|\\
			&=\frac{1}{\qty|B|}\frac{1}{\qty|g(x)|}\qty|B-g(x)|\\
			&\stackrel{\eqref{eq:sats:div-g-v.:1}}{<}\frac{1}{\qty|B|}\frac{2}{\qty|B|}\qty|B-g(x)|\\
			&\stackrel{\eqref{eq:sats:div-g-v.:2}}{<}\frac{2}{\qty|B|^2}\frac{\qty|B|^2}{2}\eps\\
			&=\eps
		\end{align*}
		Resten av satsen följer nu av Sats~\ref{sats:mult-g-v.}.
	\end{proof}
\end{theorem}

\subsection{Linjär algebra}

\subsubsection{Baser}

\begin{theorem}[Bassatsen]
	Satsen lyder, i flera delar, att:
	\begin{enumerate}
		\item \begin{enumerate}[(a)]
			\item Fler än $n$ vektorer i $\mathbb{R}^n$ är linjärt beroende.
			\item Det krävs minst $n$ vektorer för att spänna upp $\mathbb{R}^n$
			\item En bas för $\mathbb{R}^n$ har exakt $n$ vektorer.
		\end{enumerate}
		\item $\vec{v}_1,\dots,\vec{v}_n$ bas för $\mathbb{R}^n
		\iff\vec{v}_1,\dots,\vec{v}_n$ är linjärt oberoende
		$\iff\spn(\vec{v}_1,\dots,\vec{v}_n)=\mathbb{R}^n$.
	\end{enumerate}

	\begin{proof}
		Antag
		\begin{align*}
			\vec{v}_1&=(v_{11},v_{12},\dots,\vec{v}_{1n})\in\mathbb{R}^n\\
			\vdots\,\,&\\
			\vec{v}_p&=(v_{p1},v_{p2},\dots,\vec{v}_{pn})\in\mathbb{R}^n\\
			\vec{b}&=(b_1,b_2,\dots,b_n)\in\mathbb{R}^n
		\end{align*}
		Vi kan skriva om vektorekvationen
		\begin{equation}\label{eq:sats:bas-sts:1}
			x_1\vec{v}_1+\dots+x_n\vec{v}_n=\vec{b},
		\end{equation}
		där $x\in\mathbb{R}$ som de augmenterade matrisen:
		\begin{equation}\label{eq:sats:bas-sts:2}
			\overset{
				\begin{matrix}
					\color{gray}x_1
					&\quad\,\,
					&\color{gray}x_n
					&\,\,\,\,
				\end{matrix}
			}{
				\begin{amatrix}{3}
					v_{11}&\dots&v_{p1}&b_1\\
					\vdots&&\vdots&\vdots\\
					v_{1n}&\dots&v_{pn}&b_n
				\end{amatrix}%
			}\sim\begin{amatrix}{5}
				a&&&&&A\\
				0&b&&\text{\tiny Konst}&&B\\
				\vdots&0&\ddots&&&\vdots\\
				\vdots&\vdots&0&\ddots&&\vdots\\
				0&0&0&0&c&C\\
			\end{amatrix}
		\end{equation}
		Låt $k$ vara antalet pivotelement. Då gäller det att $(k\leq p)\wedge(k\leq n)$.
		\begin{enumerate}
			\item \begin{enumerate}[(a)]
				\item Om $p>n$ har LES:et \eqref{eq:sats:bas-sts:2} $p$ obekvanta
				för $n$ ekvationer. Detta innebär att det finns minst en frivariabel.
				Därmed finns ickétriviella lösningar, då $x_i\neq0\forall x_i$.
				Altså är vektorerna $\vec{v}_1,\dots,\vec{v}_p$ linjärt beroende
				om $p>n$.
				\begin{der}
					Detta betyder att för en bas gäller det att
					\begin{equation}\label{eq:sats:bas-sts:3}
						p\leq n.
					\end{equation}
				\end{der}
				\item Om $p<n$, är $S=\spn(\vec{v}_1,\dots,\vec{v}_p)$ är ett delrum
				till $\mathbb{R}^n$ s.a. $\dim(S)\leq p$. Eftersom $p<n$, $S\subsetneq\mathbb{R}^n$
				då $\dim(\mathbb{R}^n)=n>p\geq\dim(S)$. Altså krävs $p\geq n$ för att
				$\vec{v}_1,\dots,\vec{v}_p$ ska kunna spänna upp $\mathbb{R}^n$.
				\begin{der}
					Altså för en bas krävs det att
					\begin{equation}\label{eq:sats:bas-sts:4}
						p\geq n.
					\end{equation}
				\end{der}
				\item En bas av vektorerna $\vec{v}_1,\dots,\vec{v}_n$ för $\mathbb{R}^n$ kräver att:
				\begin{itemize}
					\item $\spn(\vec{v}_1,\dots,\vec{v}_n)=\mathbb{R}^n$.
					\item $\vec{v}_1,\dots,\vec{v}_n$ är linjärt oberoende.
				\end{itemize}
				Genom att applicera delsats 1. (a) och 1. (b) får vi då att:
				$$
					(p\leq n)\wedge(p\geq n)\iff p=n
				$$
				Altså måste en bas för $\mathbb{R}^n$ ha exakt $n$ vektorer.
 			\end{enumerate}
			\item Antag att vektorerna $\vec{v}_1,\dots,\vec{v}_n$ är $n$ stycken
			För att bevisa delsats 2. bryter vi upp den i 3 delar:
			\begin{enumerate}[(i)]
				\item Följden
				$$
					\vec{v}_1,\dots,\vec{v}_n\text{ bas för }\mathbb{R}^n\implies\vec{v}_1,\dots,\vec{v}_n\text{ linjärt oberoende}
				$$
				och
				$$
					\vec{v}_1,\dots,\vec{v}_n\text{ bas för }\mathbb{R}^n\implies\spn(\vec{v}_1,\dots,\vec{v}_n)=\mathbb{R}^n
				$$
				följer per definition av en bas.
				\item För att visa följden
				$$
					\vec{v}_1,\dots,\vec{v}_n\text{ linjärt oberoende }\implies\vec{v}_1,\dots,\vec{v}_n,
				$$ 
				antag att $\vec{v}_1,\dots,\vec{v}_n$ är linjärt oberoende, och för
				motsägelse att:
				$$
					S\leteq\spn(\vec{v}_1,\dots,\vec{v}_n)\neq\mathbb{R}^n.
				$$
				Då är $S\subsetneq\mathbb{R}^n$ s.a. $\dim(S)\leteq d<n$.

				Mängden $\qty{\vec{v}_1,dots,\vec{v}_n}$ har $n$ stycken vektorer
				i $S\cong\mathbb{R}^d$. Delsats 1. (a) ger att fler än $d$ vektorer
				i $\mathbb{R}^d$ är linjärt beroende. Eftersom $n>d$, är
				$\qty{\vec{v_1},\dots,\vec{v}_n}$ linjärt beroende $\contra$. Altså
				är $\spn(\vec{v}_1,\dots,\vec{v}_n)=\mathbb{R}^n$, och $\vec{v}_1,\dots,\vec{v}_n$
				ger en bas för $\mathbb{R}^n$.
				\item För att visa följden
				$$
					\spn(\vec{v}_1,\dots,\vec{v}_n)=\mathbb{R}^n\implies\vec{v}_1,\dots,\vec{v}_n,
				$$ 
				antag att $\spn(\vec{v}_1,\dots,\vec{v}_n)=\mathbb{R}^n$, och för 
				motsägelse att $\vec{v}_1,\dots,\vec{v}_n$ är linjärt beroende. Då
				gäller det att $\exists\vec{v}_i$ som är en linjär kombination av dem
				andra. Om $\vec{v}_i$ tas bort så spänner de återstående $n-1$ vektorerna
				fortfarande upp $\mathbb{R}^n$. Delsats 1. (b) däremot ger dock att
				$n$ vektorer krävs för att spänna upp $\mathbb{R}^n\contra$.

				Altså är $\vec{v}_1,\dots,\vec{v}_n$ linjärt oberoende och därmed
				ger en bas för $\mathbb{R}^n$.
			\end{enumerate}
		\end{enumerate}
	\end{proof}
\end{theorem}

\begin{lemma}
	$\vec{e}_1,\dots,\vec{e}_n$ är en bas för $\mathbb{R}^n\iff\forall\vec{v}\in\mathbb{R}^n
	\exists!v_1,\dots,v_n\in\mathbb{R}^n$ s.a.:
	$$
		\vec{v}=v_1\vec{e}_1+\dots+v_n\vec{e}_n
	$$

	\begin{proof}
		Om $\vec{e}_1,\dots,\vec{e}_n$ är en bas för $\mathbb{R}^n$, så gäller det att:
		\begin{align*}
			&\spn(\vec{e}_1,\dots,\vec{e}_n)=\mathbb{R}^n\\
			\iff&\forall\vec{v}\in\mathbb{R}^n\exists v_1,\dots,v_n\SA\vec{v}=v_1\vec{e}_1+\dots+v_n\vec{e}_n
		\end{align*}.
		Detta visar dess existens. Antag seadn att det finns två unika uppsättningar
		av konstanter som beskriver $\vec{v}$ i basen:
		\begin{align*}
			\vec{v}&=v_1\vec{e}_1+\dots+v_n\vec{e}_n\\
			\vec{v}&=v_1^*\vec{e}_1+\dots+v_n^*\vec{e}_n
		\end{align*}
		Betrakta sedan dess skillnad:
		\begin{align*}
			\vec{v}-\vec{v}=\vec{0}
			&=(v_1\vec{e}_1+\dots+v_n\vec{e}_n)-(v_1^*\vec{e}_1+\dots+v_n^*\vec{e}_n)\\
			&=(v_1-v_1^*)\vec{e}_1+\dots+(v_n-v_n^*)\vec{e}_n
		\end{align*}
		Per definition av att $\vec{e}_1,\dots,\vec{e}_n$ är en bas, så är de
		linjärt oberoende. Detta innebär att den enda lösningen till ovanstående
		ekvationen är den triviella lösningen då alla koefficienter är $0$. Altså:
		\begin{align*}
			v_1-v_1^*=0&\iff v_1=v_1^*\\
			&\quad\,\,\vdots\\
			v_n-v_n^*=0&\iff v_n=v_n^*\quad\contra
		\end{align*}
		Altså, $\forall\vec{v}\in\mathbb{R}^n\exists!v_1,\dots,v_n\in\mathbb{R}^n$ s.a.:
		$$
			\vec{v}=v_1\vec{e}_1+\dots+v_n\vec{e}_n
		$$
	\end{proof}
\end{lemma}